\documentclass[11pt,oneside]{report}

\usepackage[margin=1in]{geometry}
\usepackage{amsmath,amssymb}
\usepackage{booktabs}
\usepackage{enumitem}
\usepackage[hidelinks]{hyperref}
\usepackage[T1]{fontenc}
\usepackage{microtype}
\usepackage{xcolor}

\definecolor{codegray}{RGB}{245,245,245}
\newcommand{\code}[1]{\texttt{#1}}
\newcommand{\shape}[1]{\ensuremath{\mathtt{#1}}}
\newcommand{\btheta}{\boldsymbol{\theta}}
\newcommand{\bd}{\mathbf{d}}

\title{How the CoCoA Emulator Code Works\\
       \large A tensor-by-tensor guide from generated cosmologies to inference}
\author{V. Miranda}
\date{Working draft: Chapter 1}

\begin{document}

\pagenumbering{roman}
\maketitle
\tableofcontents
\clearpage
\pagenumbering{arabic}

\chapter{From cosmologies and data vectors to a staged training source}
\label{chap:staging}

Suppose the data-generation calculation has finished.  We have a table of
cosmologies, a list that names the cosmological parameters, and one data
vector for each cosmology.  What does the emulator code do first?

There are two answers because the program has a control path and a numerical
path.  The control path first reads the YAML file, resolves its file names,
checks which emulator family was requested, and chooses the corresponding
model class.  No cosmological row has been transformed yet.  The first
numerical operation is \emph{staging}: the code aligns the parameter table
with the data-vector dump, selects a reproducible subset of valid rows, and
decides whether that subset lives in ordinary memory or remains backed by a
file on disk.

This chapter follows that process in execution order.  It stops just before
the code constructs the whitening geometries.  Chapter~2 will begin with the
staged arrays and derive the parameter-space and data-vector transforms.

\section{The four files that must agree}

For an ordinary data-vector emulator, one training source is described by
four files.  The validation source has the same structure but different
rows.

\begin{center}
\begin{tabular}{p{0.24\textwidth}p{0.29\textwidth}p{0.37\textwidth}}
\toprule
File & Contents & Why the code needs it \\
\midrule
Parameter table, \code{*.txt}
  & One row per cosmology
  & Supplies the input parameters $\btheta_i$. \\
Parameter names, \code{*.paramnames}
  & One name per parameter-table column
  & States which physical parameter occupies each column. \\
Parameter covariance, \code{*.covmat}
  & A square covariance matrix with a name header
  & Supplies a second declaration of the input order and later defines the
    input whitening transform. \\
Data-vector dump, \code{*.npy}
  & One vector $\bd_i$ per cosmology
  & Supplies the targets that the network will learn. \\
\bottomrule
\end{tabular}
\end{center}

Let $N$ be the number of generated cosmologies, $P$ the number of modeled
parameters, and $D$ the length of one physical data vector.  The two arrays
that matter numerically are
\begin{equation}
  C =
  \begin{bmatrix}
    \btheta_0^{\mathsf T} \\
    \btheta_1^{\mathsf T} \\
    \vdots \\
    \btheta_{N-1}^{\mathsf T}
  \end{bmatrix}
  \in \mathbb{R}^{N\times P},
  \qquad
  V =
  \begin{bmatrix}
    \bd_0^{\mathsf T} \\
    \bd_1^{\mathsf T} \\
    \vdots \\
    \bd_{N-1}^{\mathsf T}
  \end{bmatrix}
  \in \mathbb{R}^{N\times D}.
  \label{eq:raw-arrays}
\end{equation}
Row $i$ of $C$ and row $i$ of $V$ must describe the same cosmology.  This
row identity is the central invariant of staging.  A network trained on
$C_i$ paired with $V_j$ for $i\ne j$ does not learn a noisy approximation.
It learns the wrong function.

\paragraph{The current parameter-table layout.}
The non-scalar loader currently reads the complete text table and then uses
the positional slice \code{[:, 2:-1]}.  In other words, it assumes that the
first two columns are the sample weight and minus log posterior, the modeled
parameters occupy the middle columns, and the last column is not a model
input.  The \code{.paramnames} file is checked against the covariance header
when it is present, but the returned name resolution is not yet used to form
the slice.  This is a current implementation constraint, not a property of
the mathematics.  A table with extra derived columns in the middle can pass
the name comparison and still violate the positional assumption.

The scalar-family loader is different.  Its inputs and targets are both
columns of the parameter table, so it resolves the requested input and output
names explicitly through the \code{.paramnames} sidecar.

\section{Control path: resolving a run before touching the rows}

The single-run driver is \code{cosmic\_shear\_train\_emulator.py}.  Its
family siblings call the same package layer.  The startup chain is
\begin{equation*}
\begin{array}{c}
\text{command-line strings}\\
\downarrow\quad \code{resolve\_cocoa\_config}\\
\text{parsed YAML with absolute data paths}\\
\downarrow\quad \code{EmulatorExperiment.from\_config}\\
\text{validated family, model class, device, and run settings}\\
\downarrow\quad \code{exp.run()}\\
\text{stage train, stage validation, build geometries, train.}
\end{array}
\end{equation*}

The path resolver reads the environment variable \code{ROOTDIR}.  It joins
the project name and the \code{chains} directory to each relative data-file
name.  An absolute path remains absolute.  This step changes strings in a
Python dictionary.  It does not load the large arrays.

\code{EmulatorExperiment.from\_config} is an alternative constructor.  A
\code{classmethod} receives the class as its first argument, conventionally
named \code{cls}, and returns a newly constructed instance.  It validates the
configuration and chooses one branch: cosmic shear, scalar outputs, a CMB
spectrum, a one-dimensional background grid, or a two-dimensional
matter-power grid.  The ordinary constructor then stores cheap state:
parameter names, the selected compute device, the logger, and empty slots for
objects that do not exist yet.

The important empty slots are
\begin{center}
\begin{tabular}{p{0.26\textwidth}p{0.63\textwidth}}
\toprule
Attribute & Meaning before and after the corresponding method runs \\
\midrule
\code{train\_set}
  & \code{None} initially; a staged training-source dictionary after
    \code{stage\_train}. \\
\code{val\_set}
  & \code{None} initially; a staged validation-source dictionary after
    \code{stage\_val}. \\
\code{pgeom}
  & \code{None} initially; the parameter transform after
    \code{build\_geometry}. \\
\code{geom}
  & \code{None} initially; the output transform after
    \code{build\_geometry}. \\
\code{chi2fn}
  & \code{None} initially; the loss wrapper after
    \code{build\_geometry}. \\
\code{model}
  & \code{None} initially; the trained PyTorch module after
    \code{train}. \\
\bottomrule
\end{tabular}
\end{center}

This delayed construction matters in parameter sweeps.  A sweep can keep the
validation source fixed while replacing the training subset, or keep both
sources fixed while rebuilding only a model.  Expensive state is therefore
created by named methods instead of being hidden inside the constructor.

\section{Numerical path: opening the source}

\code{stage\_train} creates a local random-number generator from the YAML
value \code{split\_seed}, then calls \code{load\_source} in
\code{emulator/data\_staging.py}.  The generator is local state.  It is passed
as an argument, so row selection does not depend on a module-level random
state left behind by an earlier call.

The two large inputs are opened differently:
\begin{verbatim}
dv = np.load(dv_path, mmap_mode="r", allow_pickle=False)
C  = np.loadtxt(params_path, dtype="float32")[:, param_cols]
\end{verbatim}

\code{np.loadtxt} is eager: it parses the complete parameter table and creates
a new in-memory NumPy array.  The requested dtype is \code{float32}, four
bytes per element.  The slice that follows selects the modeled columns.

\code{np.load(..., mmap\_mode="r")} is lazy with respect to the payload.  It
creates an \code{np.memmap}, an array-like object whose elements remain in the
\code{.npy} file until a later indexing operation asks the operating system
to read them.  The mode \code{"r"} makes this mapping read-only.  The code can
inspect \code{dv.shape} without loading $N D$ floating-point values.

At this point the data still live on the host, meaning ordinary CPU memory or
disk-backed virtual memory.  Nothing has moved to CUDA or Apple MPS.  Device
placement occurs later, after the geometry exists and the batching code knows
which rows a particular minibatch needs.

\subsection{The first row-identity check}

The loader immediately compares the row counts:
\begin{equation}
  C.\mathrm{shape}[0] = V.\mathrm{shape}[0].
\end{equation}
If the counts differ, it raises an error that names both files and both
counts.  Equality is necessary but not sufficient.  It proves that the files
have the same number of rows; it cannot prove by itself that row $i$ in one
file belongs to row $i$ in the other.  The generator must preserve that
ordering when it writes the pair.

\section{A seeded permutation, followed by physical cuts}

The loader next creates a permutation of the integers $0,\ldots,N-1$:
\begin{verbatim}
order = torch.randperm(n, generator=gen).numpy()
\end{verbatim}
\code{torch.randperm} returns a one-dimensional PyTorch tensor containing
every integer exactly once in pseudorandom order.  It is not an iterator and
it is not lazy: all $N$ indices are materialized.  The call to
\code{.numpy()} exposes the CPU tensor as a NumPy array.  Because both objects
use CPU storage, this conversion normally shares memory rather than copying
the index values.

The physical-window function receives $C$, the permutation, and the ordered
parameter names.  Each active window locates its columns by name and computes
one derived quantity per candidate row:
\begin{align}
  \omega_b h^2 &= \Omega_b\left(\frac{H_0}{100}\right)^2,\\
  (\Omega_m h)^2 &= \left(\Omega_m\frac{H_0}{100}\right)^2,\\
  \omega_m &= \Omega_m\left(\frac{H_0}{100}\right)^2,\\
  \omega_m n_s &=
  \Omega_m\left(\frac{H_0}{100}\right)^2 n_s.
\end{align}
For an active lower bound $a$ and upper bound $b$, the comparison is strict:
$a<q_i<b$.  The code builds a Boolean mask for each window and intersects it
with the masks already applied.  A Boolean mask is an array of \code{True}
and \code{False} values, one per candidate row.  Indexing
\code{order[mask]} uses NumPy advanced indexing and creates a new array of
the surviving integer indices.  Their order remains the shuffled order.

The result, called \code{phys}, is therefore not a new cosmology table.  It is
a list of row coordinates into the original files:
\begin{equation*}
  \code{phys} = [i_0,i_1,\ldots],
  \qquad C[i_k]\text{ and }V[i_k]\text{ remain a matched pair.}
\end{equation*}

\section{Selecting an absolute number of rows}

The YAML states an absolute number of training rows, \code{n\_train}, and an
absolute number of validation rows, \code{n\_val}.  After the cuts, the
loader checks that the requested number is at least one and that the physical
pool is large enough.  It then takes
\begin{verbatim}
idx = phys[:n_keep]
\end{verbatim}
The slice is the first $n_{\rm keep}$ entries of the seeded, cut permutation.
Using the same seed at multiple training sizes creates nested learning-curve
subsets: the $2{,}000$-row set is the prefix of the $4{,}000$-row set rather
than an unrelated draw.

Training and validation use different files.  Reusing the numerical seed
therefore repeats the selection algorithm, not the cosmologies themselves.
The scientific assumption is that the two file sets are disjoint.  That
assumption must be checked from file identity or row identity; a seed does not
establish disjointness.

\section{Staging: copy the subset or retain the disk mapping}

The word \emph{stage} means ``put the selected rows in the storage form from
which later batches will read them.''  It does not yet mean ``move them to the
GPU.''  The helper \code{stage\_source} estimates whether the selected rows
fit within a configured fraction of available host memory.  In the current
implementation, this estimate counts the selected data-vector array but not
the selected parameter-array copy.  The resident branch nevertheless creates
both arrays.  The estimate can therefore be low by
$n_u P\,\mathrm{sizeof}(C_{ij})$ bytes.  This is a memory-accounting defect in
the present code, not a license for consumers to ignore the parameter copy.

There are two outcomes.

\subsection{Resident host-memory outcome}

If the pair fits, the helper gathers the sorted unique selected rows into
compact arrays.  Advanced integer indexing is eager and returns a copy.  The
compact arrays have shapes
\begin{equation}
  C_{\rm src}\in\mathbb{R}^{n_u\times P},
  \qquad
  V_{\rm src}\in\mathbb{R}^{n_u\times D},
\end{equation}
where $n_u$ is the number of distinct selected rows.  The source's
\code{idx} is then expressed in the compact arrays' local row coordinates.
If original dump row 91 becomes compact row 3, later code asks the resident
array for row 3, not row 91.

\subsection{Disk-backed outcome}

If the data vectors do not fit, the helper retains the original read-only
\code{np.memmap}.  The source's \code{idx} remains in global on-disk row
coordinates.  A later batch read such as \code{dv[idx]} asks the operating
system for only those rows.  This advanced indexing operation materializes
the requested batch; it does not turn the entire dump into an in-memory
array.

The parameter table was already loaded eagerly by \code{np.loadtxt}.  Thus,
``disk-backed source'' describes the data vectors, not every object in the
source dictionary.

\section{The source dictionary and its coordinate systems}

\code{load\_source} returns a small dictionary:
\begin{verbatim}
src = {"C":         C_src,
       "dv":        dv_src,
       "idx":       idx_src,
       "dump_rows": sorted_unique_original_rows}
\end{verbatim}

The fields have distinct jobs.

\begin{center}
\begin{tabular}{p{0.18\textwidth}p{0.25\textwidth}p{0.47\textwidth}}
\toprule
Key & Coordinate system & Meaning \\
\midrule
\code{C}
  & Matches \code{idx}
  & Parameter rows.  It is compact in the resident outcome and the full
    eager table in the disk-backed outcome. \\
\code{dv}
  & Matches \code{idx}
  & Data-vector rows.  It is compact in the resident outcome and a file
    mapping in the disk-backed outcome. \\
\code{idx}
  & Local or global
  & The row coordinates that consumers must use with this particular
    \code{C}/\code{dv} pair. \\
\code{dump\_rows}
  & Always global
  & Sorted unique row numbers in the original files.  A sibling dump uses
    these coordinates to retrieve the same cosmologies. \\
\bottomrule
\end{tabular}
\end{center}

\code{dump\_rows} is needed by the two-dimensional matter-power path.  That
path has a second, row-aligned base dump.  Even when staging has converted
the primary arrays to local coordinates, the sibling file still needs the
original disk row numbers.  Sorting makes its reads sequential.  Applying
\code{np.unique} removes duplicates and also sorts; it returns a new array.

The complete coordinate flow is
\begin{equation*}
\begin{array}{c}
\text{original rows }0,\ldots,N-1\\
\downarrow\quad\text{seeded permutation and physical mask}\\
\code{idx}:\text{ selected rows in shuffled training order}\\
\downarrow\quad\text{sorted unique gather, only if copied to RAM}\\
\code{idx\_src}:\text{ local coordinates into compact arrays}\\[2mm]
\code{dump\_rows}:\text{ original coordinates retained for sibling files.}
\end{array}
\end{equation*}

\section{Training statistics are owned by the training source}

For the training source, \code{with\_means=True}.  The loader computes the
mean of every parameter column and every staged data-vector column:
\begin{equation}
  \bar C_j = \frac{1}{n}\sum_{i=1}^{n}C_{ij},
  \qquad
  \bar V_k = \frac{1}{n}\sum_{i=1}^{n}V_{ik}.
\end{equation}
These arrays are stored as \code{C\_mean} and \code{dv\_mean}.  The
validation source uses \code{with\_means=False}; it does not estimate a
second zero point.  Training and validation must be expressed in the same
coordinate system, so the training-derived geometry transforms both.

Subtracting a center does not change a residual-based loss.  If prediction
and truth are both centered by $\bar V$, then
\begin{equation}
  (V_{\rm pred}-\bar V)-(V_{\rm true}-\bar V)
  = V_{\rm pred}-V_{\rm true}.
\end{equation}
Centering changes the zero point the network learns, not the physical
prediction-minus-truth residual.

The matter-power path is again special.  Its raw surface is transformed into
``law space''---for example, $\log(P/P_{\rm base})$---and may be thinned along
the wavenumber axis.  Its final output mean must therefore be computed after
that transformation.  The generic loader skips the raw \code{dv\_mean}; the
bounded matter-power staging code streams the transformed mean and variance
over row chunks.

\section{What has not happened yet}

At the end of \code{stage\_train} and \code{stage\_val}, the program has
selected and located the rows.  It has not yet built a neural network, an
optimizer, or a covariance-weighted loss.  Apart from the seeded permutation,
the arrays are still NumPy objects on the host.

The next method, \code{build\_geometry}, creates two maps:
\begin{align*}
  \text{parameter geometry: }&
  \btheta\longmapsto x,
  &&\text{raw cosmology to centered, rescaled network input},\\
  \text{output geometry: }&
  \bd\longmapsto t,
  &&\text{physical data vector to centered, covariance-scaled target}.
\end{align*}
Only after those maps exist can the batching code convert a NumPy batch with
\code{torch.from\_numpy}, cast it to the model's \code{float32} dtype, move it
to the selected device, and encode it.  That is the subject of Chapter~2.

\section{Chapter summary}

The first numerical job is not neural-network training.  It is identity
preservation.  The staging code must keep four statements true:
\begin{enumerate}[leftmargin=2em]
  \item Parameter row $i$ and data-vector row $i$ describe the same
        cosmology.
  \item Parameter names and column order match the covariance that will
        transform them.
  \item The seeded subset is selected after the declared physical cuts and
        contains the requested absolute number of rows.
  \item Every row index is interpreted in the coordinate system of the array
        it indexes: compact local rows for a resident copy, original global
        rows for a file-backed dump.
\end{enumerate}
Once those conditions hold, the program has a staged training source.  The
next problem is numerical geometry: how to transform correlated parameters
and correlated outputs into coordinates a neural network can learn without
silently changing the $\chi^2$ that defines emulator accuracy.

\end{document}
